%% 02_related_work.tex
\section{Related Work}
\label{sec:related}

\paragraph{Clinical and medical RAG.}
Lewis et al.~\citep{lewis2020retrieval} introduced RAG as a paradigm for knowledge-intensive
NLP tasks; the MIRAGE benchmark~\citep{xiong2024benchmarking} established medical RAG evaluation
over literature corpora, and Almanac~\citep{zakka2024almanac} demonstrated clinical guideline
grounding.  These systems retrieve over \emph{population-level} knowledge (publications,
guidelines), not per-patient EHR data.  Schmiedmayer et al.~\citep{schmiedmayer2024llmonfhir}
built LLM~on~FHIR, a function-calling retrieval system over patient-level FHIR bundles,
validated on Synthea data with physician Likert ratings.  The peer-reviewed version,
LLMonFHIR~\citep{schmiedmayer2025llmonfhir}, is the closest functional predecessor to our
work; it evaluates only structured retrieval and not a narrative baseline.
MedGraphRAG~\citep{wu2024medgraphrag} and KG-RAG~\citep{soman2024biomedical} add
hierarchical graph structure to retrieval, but over LLM-extracted or literature-derived graphs
rather than schema-defined FHIR resources.

\paragraph{Structured vs.\ unstructured EHR.}
Moldwin et al.~\citep{moldwin2021empirical} performed the closest paired same-patient
comparison, finding that text features outperformed structured features for 84\% of
MIMIC phenotypes while acknowledging the information-content confound explicitly.  TabLLM~\citep{hegselmann2023tabllm}
and Lee et al.~\citep{lee2025pseudonotes} demonstrate competitive or superior performance from
serialising structured EHR tables to natural language for classification tasks.  DrugEHRQA~\citep{bardhan2022drugehrqa}
pairs structured tables and discharge notes for medication QA but admits cross-modal
information asymmetry.  CLEAR~\citep{lopez2025clinical} compares entity-aware vs.\ embedding
RAG on 20,000 clinical notes, reporting F1 gains from entity-structured retrieval; the
comparison is within-narrative, not FHIR-structured.

\paragraph{FHIR and ML/LLM.}
FHIR-Former~\citep{engelke2025fhirformer} applies FHIR-structured resources plus notes to
clinical prediction tasks.  SQL-on-FHIR~\citep{grimes2025sqlonfhir} and Pathling~\citep{grimes2022pathling}
provide analytics-oriented FHIR access, complementary to but distinct from semantic retrieval.
FHIR-GPT~\citep{li2024fhirgpt} and the FHIR Workbench~\citep{idrissiyaghir2025fhirworkbench}
characterise LLM competence on FHIR production and understanding tasks.
Kothari and Gupta~\citep{kothari2025patientqa} propose a two-stage pipeline for FHIR
QA on Synthea (the closest Synthea-grounded predecessor) but include no narrative baseline
and focus on fine-tuning rather than retrieval representation.
FHIR-AgentBench~\citep{lee2025fhiragentbench} benchmarks agentic retrieval strategies over
MIMIC-IV-FHIR at scale, identifying medication questions as the weakest category
(5--14\% correctness) and retrieval precision as the dominant bottleneck; it contains no
narrative arm.

\paragraph{Clinical temporal reasoning.}
Temporal reasoning remains the weakest category in clinical LLM performance across benchmark
evaluations. Zero-shot temporal-relation extraction substantially trails supervised methods
on clinical text~\citep{yuan2023zeroshot}, and instruction-following on longitudinal EHRs
remains hard even for frontier models: MedAlign~\citep{fleming2024medalign}, a dataset of
983 clinician-curated instructions over 275 OMOP-structured longitudinal records, shows
GPT-4 reaching only 60.1\% correctness despite a 32k context window.
TIMER~\citep{cui2025timer} introduces a dedicated temporal instruction benchmark for
longitudinal clinical records and reports a 6.6\% completeness improvement from specialised
instruction tuning, confirming that temporal grounding is a distinct skill not inherited
from general language modelling.
Our error taxonomy provides a complementary characterisation of temporal-anchor failures
specifically within structured-FHIR and narrative RAG retrieval contexts.

\paragraph{Adherence-indicator ML.}
Lo-Ciganic et al.~\citep{lociganic2015using} established the canonical feature set for
claims-based adherence prediction; Zullig et al.~\citep{zullig2019novel} provide a
reference AUC benchmark (0.664--0.677).  Nateghi Haredasht et al.~\citep{nateghi2025predicting}
demonstrate LLM-extracted narrative features for medication retention prediction on a MOUD
cohort; Anderson et al.~\citep{anderson2025paging} apply the same design to mortality.
Gao et al.~\citep{gao2024raw} show that raw numerical features often match or beat LLM
embeddings on tabular clinical tasks.  Our O6 arm extends this line to a non-MOUD
specialty-regimen adherence task with matched FHIR-structured and narrative feature sets.

\paragraph{Novelty claim.}
To our knowledge, no prior work evaluates narrative RAG, naive structured FHIR RAG, and
resource-aware structured FHIR RAG on \emph{paired data} for specialty-medication
adherence-indicator questions with programmatic ground truth, nor compares these
representations as feature sources for adherence classification on the same patient cohort.
The four-part intersection -- (1) paired same-bundle rendering, (2) resource-aware vs.\
naive structured retrieval, (3) specialty-medication QA domain, (4) feature-extraction
comparison -- is unoccupied in the literature we surveyed.
